% GENERATED FILE. Edit canonical lessons, then run npm run generate:books.
% canonical-source-hash: fnv1a64:021e06eeade0198b
% canonical-lessons: HI-C537-karyakram, HI-C537-mohlat, HI-C537-harbari, HI-C537-muhurt, HI-C537-khiskana, HI-R537-review-530-537

\chapter{A Change of Programme}
\label{ch:hi-a-change-of-programme}
\hlchaptermodality{537}

\begin{chapteropening}
\textbf{By the end of this chapter:} \emph{I can say that something was moved, rushed or given more time.}

Complete the last lesson of chapter 537: I can say that something was moved, rushed or given more time.
\end{chapteropening}

\section[\texorpdfstring{\hi{कार्यक्रम} (kāryakram)}{कार्यक्रम (kāryakram)}]{\hi{कार्यक्रम} (kāryakram) — a programme; a schedule}

\label{lesson:HI-C537-karyakram}



\begin{quote}
Before the new one: say the Hindi for an outline, then the Hindi for a draft (of a text).
\end{quote}

\subsection*{You'll want to know: \hi{कार्यक्रम}}
\textbf{\hi{कार्यक्रम}} — \emph{kāryakram} — \textquotedblleft{}a programme; a schedule\textquotedblright{}.

\begin{quote}
\hi{स्कूल} \hi{का} \hi{कार्यक्रम} \hi{शाम} \hi{छह} \hi{बजे} \hi{है।}
\end{quote}

\begin{quote}
\emph{Skūl kā kāryakram śām chah baje hai.}
\end{quote}

\begin{quote}
The school programme is at six in the evening.
\end{quote}

Masculine.

\begin{cousinweb}
\textbf{\hi{कार्यक्रम}} (\emph{kāryakram}) is a modern compound of two Sanskrit words, \emph{kārya}, \textquotedblleft{}to be done\textquotedblright{}, and \emph{krama}, \textquotedblleft{}a step, a course\textquotedblright{} (Monier-Williams): the things to be done, one after another.
\end{cousinweb}

\begin{grammarlens}[title={What you've built}]
One more word for saying what happens if.
\end{grammarlens}

\subsection*{Your turn}
\begin{itemize}
\raggedright
  \item \emph{Say it:} \emph{kāryakram}
  \item \emph{Say it:} \emph{kāryakram}, once more
  \item \emph{Say it:} \emph{Skūl kā kāryakram śām chah baje hai.} — The school programme is at six in the evening.
  \item \emph{Recall:} say the Hindi for an outline, then the Hindi for a draft (of a text), then say \emph{kāryakram} again
\end{itemize}

\begin{tcolorbox}[breakable,colback=teal!4,colframe=teal!35!black,title={Before you move on}]
What does \hi{कार्यक्रम} mean? (\textquotedblleft{}A programme; a schedule\textquotedblright{}.) Say it once more.
\end{tcolorbox}
\section[\texorpdfstring{\hi{मोहलत} (mohlat)}{मोहलत (mohlat)}]{\hi{मोहलत} (mohlat) — extra time; a respite}

\label{lesson:HI-C537-mohlat}



\begin{quote}
Before the new one: say the Hindi for a draft (of a text), then the Hindi for a programme.
\end{quote}

\subsection*{You'll want to know: \hi{मोहलत}}
\textbf{\hi{मोहलत}} — \emph{mohlat} — \textquotedblleft{}extra time; a respite\textquotedblright{}.

\begin{quote}
\hi{मालिक} \hi{ने} \hi{हमें} \hi{किराया} \hi{देने} \hi{के} \hi{लिए} \hi{एक} \hi{हफ़्ते} \hi{की} \hi{मोहलत} \hi{दी।}
\end{quote}

\begin{quote}
\emph{Mālik ne hamẽ kirāyā dene ke liye ek hafte kī mohlat dī.}
\end{quote}

\begin{quote}
The landlord gave us a week's grace to pay the rent.
\end{quote}

Feminine.

\begin{cousinweb}
\textbf{\hi{मोहलत}} (\emph{mohlat}) is usually traced to Arabic \emph{muhla}, \textquotedblleft{}respite, delay\textquotedblright{}, through Persian.
\end{cousinweb}

\begin{grammarlens}[title={What you've built}]
One more word for saying what happens if.
\end{grammarlens}

\subsection*{Your turn}
\begin{itemize}
\raggedright
  \item \emph{Say it:} \emph{mohlat}
  \item \emph{Say it:} \emph{mohlat}, once more
  \item \emph{Say it:} \emph{Mālik ne hamẽ kirāyā dene ke liye ek hafte kī mohlat dī.} — The landlord gave us a week's grace to pay the rent.
  \item \emph{Recall:} say the Hindi for a draft (of a text), then the Hindi for a programme, then say \emph{mohlat} again
\end{itemize}

\begin{tcolorbox}[breakable,colback=teal!4,colframe=teal!35!black,title={Before you move on}]
What does \hi{मोहलत} mean? (\textquotedblleft{}Extra time; a respite\textquotedblright{}.) Say it once more.
\end{tcolorbox}
\section[\texorpdfstring{\hi{हड़बड़ी} (haṛbaṛī)}{हड़बड़ी (haṛbaṛī)}]{\hi{हड़बड़ी} (haṛbaṛī) — a flustered hurry}

\label{lesson:HI-C537-harbari}



\begin{quote}
Before the new one: say the Hindi for a programme, then the Hindi for extra time.
\end{quote}

\subsection*{You'll want to know: \hi{हड़बड़ी}}
\textbf{\hi{हड़बड़ी}} — \emph{haṛbaṛī} — \textquotedblleft{}a flustered hurry\textquotedblright{}.

\begin{quote}
\hi{हड़बड़ी} \hi{में} \hi{मैं} \hi{अपना} \hi{बटुआ} \hi{घर} \hi{पर} \hi{भूल} \hi{गया।}
\end{quote}

\begin{quote}
\emph{Haṛbaṛī meṁ maiṁ apnā baṭuā ghar par bhūl gayā.}
\end{quote}

\begin{quote}
In my hurry I forgot my wallet at home.
\end{quote}

Feminine.

\begin{cousinweb}
\textbf{\hi{हड़बड़ी}} (\emph{haṛbaṛī}) is a tadbhava, inherited from \emph{*haḍabaḍa}, \textquotedblleft{}confusion\textquotedblright{}, a form CDIAL reconstructs from its descendants (CDIAL 13949). Today it is the confusion of someone short of time.
\end{cousinweb}

\begin{grammarlens}[title={What you've built}]
One more word for saying what happens if.
\end{grammarlens}

\subsection*{Your turn}
\begin{itemize}
\raggedright
  \item \emph{Say it:} \emph{haṛbaṛī}
  \item \emph{Say it:} \emph{haṛbaṛī}, once more
  \item \emph{Say it:} \emph{Haṛbaṛī meṁ maiṁ apnā baṭuā ghar par bhūl gayā.} — In my hurry I forgot my wallet at home.
  \item \emph{Recall:} say the Hindi for a programme, then the Hindi for extra time, then say \emph{haṛbaṛī} again
\end{itemize}

\begin{tcolorbox}[breakable,colback=teal!4,colframe=teal!35!black,title={Before you move on}]
What does \hi{हड़बड़ी} mean? (\textquotedblleft{}A flustered hurry\textquotedblright{}.) Say it once more.
\end{tcolorbox}
\section[\texorpdfstring{\hi{मुहूर्त} (muhūrt)}{मुहूर्त (muhūrt)}]{\hi{मुहूर्त} (muhūrt) — an auspicious moment (to begin something)}

\label{lesson:HI-C537-muhurt}



\begin{quote}
Before the new one: say the Hindi for extra time, then the Hindi for a flustered hurry.
\end{quote}

\subsection*{You'll want to know: \hi{मुहूर्त}}
\textbf{\hi{मुहूर्त}} — \emph{muhūrt} — \textquotedblleft{}an auspicious moment (to begin something)\textquotedblright{}.

\begin{quote}
\hi{पुजारी} \hi{ने} \hi{शादी} \hi{का} \hi{मुहूर्त} \hi{निकाला।}
\end{quote}

\begin{quote}
\emph{Pujārī ne śādī kā muhūrt nikālā.}
\end{quote}

\begin{quote}
The priest worked out the auspicious time for the wedding.
\end{quote}

Masculine.

\begin{cousinweb}
\textbf{\hi{मुहूर्त}} (\emph{muhūrt}) is a tatsama, taken whole from Sanskrit \emph{muhūrta}, \textquotedblleft{}a moment, a short space of time\textquotedblright{} (Monier-Williams). Hindi uses it for the moment the stars favour.
\end{cousinweb}

\begin{grammarlens}[title={What you've built}]
One more word for saying what happens if.
\end{grammarlens}

\subsection*{Your turn}
\begin{itemize}
\raggedright
  \item \emph{Say it:} \emph{muhūrt}
  \item \emph{Say it:} \emph{muhūrt}, once more
  \item \emph{Say it:} \emph{Pujārī ne śādī kā muhūrt nikālā.} — The priest worked out the auspicious time for the wedding.
  \item \emph{Recall:} say the Hindi for extra time, then the Hindi for a flustered hurry, then say \emph{muhūrt} again
\end{itemize}

\begin{tcolorbox}[breakable,colback=teal!4,colframe=teal!35!black,title={Before you move on}]
What does \hi{मुहूर्त} mean? (\textquotedblleft{}An auspicious moment (to begin something)\textquotedblright{}.) Say it once more.
\end{tcolorbox}
\section[\texorpdfstring{\hi{खिसकाना} (khiskānā)}{खिसकाना (khiskānā)}]{\hi{खिसकाना} (khiskānā) — to shift; to push back (a date)}

\label{lesson:HI-C537-khiskana}



\begin{quote}
Before the new one: say the Hindi for a flustered hurry, then the Hindi for an auspicious moment (to begin something).
\end{quote}

\subsection*{You'll want to know: \hi{खिसकाना}}
\textbf{\hi{खिसकाना}} — \emph{khiskānā} — \textquotedblleft{}to shift; to push back (a date)\textquotedblright{}.

\begin{quote}
\hi{हमने} \hi{बैठक} \hi{अगले} \hi{हफ़्ते} \hi{तक} \hi{खिसका} \hi{दी।}
\end{quote}

\begin{quote}
\emph{Hamne baiṭhak agle hafte tak khiskā dī.}
\end{quote}

\begin{quote}
We pushed the meeting back to next week.
\end{quote}

Past, with \hi{ने} (he or she …): \hi{उसने} \hi{मेज़} \hi{खिसकाई}, \emph{usne mez khiskāī}.

\begin{cousinweb}
\textbf{\hi{खिसकाना}} (\emph{khiskānā}) is the causative of \emph{khisaknā}, \textquotedblleft{}to slip away\textquotedblright{}, which CDIAL files under \emph{*khis}, a reconstructed form (CDIAL 3888).
\end{cousinweb}

\begin{grammarlens}[title={What you've built}]
One more word for saying what happens if.
\end{grammarlens}

\subsection*{Your turn}
\begin{itemize}
\raggedright
  \item \emph{Say it:} \emph{khiskānā}
  \item \emph{Say it:} \emph{khiskānā}, once more
  \item \emph{Say it:} \emph{Hamne baiṭhak agle hafte tak khiskā dī.} — We pushed the meeting back to next week.
  \item \emph{Recall:} say the Hindi for a flustered hurry, then the Hindi for an auspicious moment (to begin something), then say \emph{khiskānā} again
\end{itemize}

\begin{tcolorbox}[breakable,colback=teal!4,colframe=teal!35!black,title={Before you move on}]
What does \hi{खिसकाना} mean? (\textquotedblleft{}To shift; to push back (a date)\textquotedblright{}.) Say it once more.
\end{tcolorbox}
\section[(dialogue)]{Looking back: from Learning a Raga to A Change of Programme}

\label{lesson:HI-R537-review-530-537}



\begin{quote}
Say the Hindi for music and a musical note.
\end{quote}

\subsection*{Your turn}
Say these aloud:

\begin{itemize}
\raggedright
  \item music, a musical note, a beat, a raga, to hum
  \item a small two-headed drum, a shehnai (a reed pipe), ankle bells, a musical instrument, a mridang (a barrel drum)
  \item a lyric, tempo, practice (of music), a gathering (for music or poetry), to make dance
  \item wrestling, a wrestler, a wrestling ground, to pit against, a stance
  \item a match (a game), a contest, to bring victory to, a winner, a prize, an award
  \item a round, a phase, quickness, to leap over, a medal, to make run
  \item a plan, a decision (formal), an arrangement, a system, an outline, a draft (of a text)
  \item a programme, extra time, a flustered hurry, an auspicious moment (to begin something), to shift
\end{itemize}

\begin{tcolorbox}[breakable,colback=teal!4,colframe=teal!35!black,title={Before you move on}]
Music? (\textbf{\hi{संगीत}}, \emph{saṅgīt}.) A musical note? (\textbf{\hi{सुर}}, \emph{sur}.) A beat? (\textbf{\hi{ताल}}, \emph{tāl}.) A raga? (\textbf{\hi{राग}}, \emph{rāg}.) To hum? (\textbf{\hi{गुनगुनाना}}, \emph{gungunānā}.) A small two-headed drum? (\textbf{\hi{ढोलक}}, \emph{ḍholak}.) A shehnai (a reed pipe)? (\textbf{\hi{शहनाई}}, \emph{śahnāī}.) Ankle bells? (\textbf{\hi{घुँघरू}}, \emph{ghũghrū}.) A musical instrument? (\textbf{\hi{बाजा}}, \emph{bājā}.) A mridang (a barrel drum)? (\textbf{\hi{मृदंग}}, \emph{mṛdaṅg}.) A lyric? (\textbf{\hi{गीत}}, \emph{gīt}.) Tempo? (\textbf{\hi{लय}}, \emph{lay}.) Practice (of music)? (\textbf{\hi{रियाज़}}, \emph{riyāz}.) A gathering (for music or poetry)? (\textbf{\hi{महफ़िल}}, \emph{mahfil}.) To make dance? (\textbf{\hi{नचाना}}, \emph{nacānā}.) Wrestling? (\textbf{\hi{कुश्ती}}, \emph{kuśtī}.) A wrestler? (\textbf{\hi{पहलवान}}, \emph{pahalvān}.) A wrestling ground? (\textbf{\hi{अखाड़ा}}, \emph{akhāṛā}.) To pit against? (\textbf{\hi{लड़ाना}}, \emph{laṛānā}.) A stance? (\textbf{\hi{पैंतरा}}, \emph{paĩtrā}.) A match (a game)? (\textbf{\hi{मैच}}, \emph{maic}.) A contest? (\textbf{\hi{मुक़ाबला}}, \emph{muqāblā}.) To bring victory to? (\textbf{\hi{जिताना}}, \emph{jitānā}.) A winner? (\textbf{\hi{विजेता}}, \emph{vijetā}.) A prize, an award? (\textbf{\hi{पुरस्कार}}, \emph{puraskār}.) A round, a phase? (\textbf{\hi{दौर}}, \emph{daur}.) Quickness? (\textbf{\hi{फुर्ती}}, \emph{phurtī}.) To leap over? (\textbf{\hi{फाँदना}}, \emph{phā̃dnā}.) A medal? (\textbf{\hi{मेडल}}, \emph{meḍal}.) To make run? (\textbf{\hi{दौड़ाना}}, \emph{dauṛānā}.) A plan? (\textbf{\hi{योजना}}, \emph{yojnā}.) A decision (formal)? (\textbf{\hi{निर्णय}}, \emph{nirṇay}.) An arrangement, a system? (\textbf{\hi{व्यवस्था}}, \emph{vyavasthā}.) An outline? (\textbf{\hi{रूपरेखा}}, \emph{rūprekhā}.) A draft (of a text)? (\textbf{\hi{मसौदा}}, \emph{masaudā}.) A programme? (\textbf{\hi{कार्यक्रम}}, \emph{kāryakram}.) Extra time? (\textbf{\hi{मोहलत}}, \emph{mohlat}.) A flustered hurry? (\textbf{\hi{हड़बड़ी}}, \emph{haṛbaṛī}.) An auspicious moment (to begin something)? (\textbf{\hi{मुहूर्त}}, \emph{muhūrt}.) To shift? (\textbf{\hi{खिसकाना}}, \emph{khiskānā}.)
\end{tcolorbox}
